\begin{table}[h!t!b!]
    \centering
    \renewcommand{\arraystretch}{1.3}
    \resizebox{\textwidth}{!}{%
    \begin{NiceTabular}{cc|ccccc|ccccc|cccc}
         & & $F^3$ & $\phi^2 F^2$ & $\psi^2\phi F$ & $\psi^4$ & $\psi^2\phi^3$ & $\bar F^3$ & $\phi^2\bar F^2$ & $\bar\psi^2\phi\bar F$ & $\bar\psi^4$ & $\bar\psi^2\phi^3$ & $\bar\psi^2\psi^2$ & $D\bar\psi\psi\phi^2$ & $D^2\phi^4$ & $\phi^6$ \\
         & $(w,\wb)$ & $(0,6)$ & $(2,6)$ & $(2,6)$ & $(2,6)$ & $(4,6)$ & $(6,0)$ & $(6,2)$ & $(6,2)$ & $(6,2)$ & $(6,4)$ & $(4,4)$ & $(4,4)$ & $(4,4)$ & $(6,6)$ \\ \hline
%        $F^3$ & $(0,6)$ & & & & & $\mathsf{L}$ & $\mathsf{H}^\star$ & \ccol$\mathsf{H}$ & \ccol$\mathsf{H}$ & \ccol$\mathsf{H}$ & $\mathsf{L}$ & \ccol$y^2$ & \ccol$y^2$ & \ccol$y^2$ & $\mathsf{L}$ \\
        $F^3$ & $(0,6)$ & & \ccol & \ccol & \ccol$\times$ & $\mathsf{L}$ & $\mathsf{H}^\star$ & \ccol$\mathsf{H}$ & \ccol$\mathsf{H}$ & \ccol$\mathsf{H}$ & $\mathsf{L}$ & \ccol$\mathsf{H}$ & \ccol$\mathsf{H}$ & \ccol$\mathsf{H}$ & $\mathsf{L}$ \\
        $\phi^2 F^2$ & $(2,6)$ & & & & & \ccol & & & & & \ccol$\mathsf{H}$ & & & & $\mathsf{L}$ \\
        $\psi^2\phi F$ & $(2,6)$ & & & & & \ccol & & & & & \ccol$\mathsf{H}$ & & & & $\mathsf{L}$ \\
        $\psi^4$ & $(2,6)$ & & & & & \ccol & & & & & \ccol$\mathsf{H}$ & & & & $\mathsf{L}$ \\
        $\psi^2\phi^3$ & $(4,6)$ & & & & & & & & & & & & & & \ccol \\ \hline
%        $\bar F^3$ & $(6,0)$ & $\mathsf{H}^\star$ & \ccol$\mathsf{H}$ & \ccol$\mathsf{H}$ & \ccol$\mathsf{H}$ & $\mathsf{L}$ & & & & & $\mathsf{L}$ & \ccol$\bar y^2$ & \ccol$\bar y^2$ & \ccol$\bar y^2$ & $\mathsf{L}$ \\
        $\bar F^3$ & $(6,0)$ & $\mathsf{H}^\star$ & \ccol$\mathsf{H}$ & \ccol$\mathsf{H}$ & \ccol$\mathsf{H}$ & $\mathsf{L}$ & & \ccol & \ccol & \ccol$\times$ & $\mathsf{L}$ & \ccol$\mathsf{H}$ & \ccol$\mathsf{H}$ & \ccol$\mathsf{H}$ & $\mathsf{L}$ \\
        $\phi^2\bar F^2$ & $(6,2)$ & & & & & \ccol$\mathsf{H}$ & & & & & \ccol & & & & $\mathsf{L}$ \\
        $\bar\psi^2\phi\bar F$ & $(6,2)$ & & & & & \ccol$\mathsf{H}$ & & & & & \ccol & & & & $\mathsf{L}$ \\
        $\bar\psi^4$ & $(6,2)$ & & & & & \ccol$\mathsf{H}$ & & & & & \ccol & & & & $\mathsf{L}$ \\
        $\bar\psi^2\phi^3$ & $(6,4)$ & & & & & & & & & & & & & & \ccol \\ \hline
        $\bar\psi^2\psi^2$ & $(4,4)$ & & & & & \ccol & & & & & \ccol & & & & $\mathsf{L}$ \\
        $D\bar\psi\psi\phi^2$ & $(4,4)$ & & & & & \ccol & & & & & \ccol & & & & $\mathsf{L}$ \\
        $D^2\phi^4$ & $(4,4)$ & & & & & \ccol & & & & & \ccol & & & & $\mathsf{L}$ \\
        $\phi^6$ & $(6,6)$ & & & & & & & & & & & & & & \\
    \end{NiceTabular}
    }
    \caption{Structure of the two-loop anomalous-dimension matrix for dimension-six operators. The entry in row $\mathcal O_i$ and column $\mathcal O_j$ corresponds to $\gamma^{(2)}_{ij}$. Entries labeled $\mathsf{L}$ vanish by the length selection rule \cite{Bern:2019wie}. The shaded cells define the scheme-independent band $\ell_i=\ell_j-1$, which is completely determined by three-particle cuts of tree-level objects. Within this band, $\mathsf{H}$ denotes entries that vanish by the helicity selection rule of \Eq{eq:dl1}. When the conditions of \Eq{eq:dl1} are not satisfied, $\times$ indicates entries for which no three-particle cuts exist. Finally, entries labeled $\mathsf{H}^\star$ vanish in the holomorphic scheme, as dictated by \Eq{eq:dl2}.}
    \label{tab:dim6}
\end{table}